\section{Introduction}

Multi-omics cancer models often concatenate heterogeneous molecular measurements, including transcriptomic, copy-number, methylation, proteomic, and mutation features. Established integration approaches include latent-variable and clustering models such as iCluster and MOFA, supervised multi-block approaches such as DIABLO, and more recent deep or graph-based fusion models such as MOGONET and multimodal survival networks \cite{shen_icluster_2009,argelaguet_mofa_2018,singh_diablo_2019,wang_mogonet_2021,valesilva_multisurv_2021}. Recent benchmarking work has also shown that deep multi-omics fusion methods must be compared against strong classical and simple neural baselines rather than evaluated as isolated models \cite{leng_benchmark_2022}. Although early-fusion models are simple and strong, they do not explicitly represent the stepwise integration of different biological layers. Liquid neural networks and closed-form continuous-time recurrent models provide an alternative: their hidden states evolve through input-dependent transition rates, allowing each input view to update a latent state before prediction \cite{hasani_ltc_2021,lechner_cfc_2022}.

The motivation for applying a Liquid/CfC model to multi-omics data is conceptually different from applying it to ordinary time series. A breast tumor sample is usually measured once, but each molecular layer provides a different view of the same biological state. A recurrent state-update model can therefore be used as a structured fusion module: copy-number, mutation, and expression features can update a shared latent representation before final subtype classification. This interpretation is attractive, but it must be tested carefully because an ordered omics sequence is not a true temporal process.

In this study, we asked whether a Liquid/CfC-style architecture is useful for cross-sectional cancer multi-omics prediction under realistic external validation. Breast cancer molecular subtype prediction was selected as the main task because subtype labels are biologically meaningful, externally available across cohorts, and strongly connected to expression programs. The prediction target was a five-class PAM50-like subtype label:
\[
y_i \in \{\mathrm{Basal}, \mathrm{Her2}, \mathrm{LumA}, \mathrm{LumB}, \mathrm{Normal}\}.
\]

We deliberately framed the analysis as a systematic evaluation rather than a single-model performance claim. First, we compared Liquid/CfC against strong linear, tree-based, boosting, and neural baselines. Second, we tested whether adding METABRIC to the training pool mitigated the small-sample limitation of TCGA-only neural modeling. Third, we introduced a reduced-capacity small-Liquid/CfC model to examine whether lowering parameter count improved external generalization. Fourth, we evaluated transfer to independent cohorts, including CPTAC 2020 for multimodal validation and SMC 2018 and SCAN-B for expression-marker validation. Fifth, after reviewer-facing sample-alignment audit, we added a core-aligned multi-cancer sensitivity analysis that excludes the sparsely covered RPPA block.

The central question was therefore not whether Liquid/CfC universally outperforms all baselines, but whether it provides a competitive and biologically plausible cross-omics fusion mechanism under rigorous multi-cohort validation.
